% !TEX root = ../Report.tex
% ======================================================================
\chapter{Evaluation}
% ======================================================================

This chapter reports the evaluation conducted on Evaluation Day (Wednesday 25th
February 2026) and via a follow-up survey. It presents both objective user
acceptance test results and subjective survey findings, and interprets them
against the project's success criteria.

\section{Evaluation Day Setup}
Evaluation Day took place in the Department of Computer Science, Room CS0.03,
Table B01, slot 13:05--13:45, with the project titled ``GiTrip --- Git for
Collaborative Trip Planning''. The demonstration environment used four Linux
machines plus the author's laptop so that multiple participants could test
branching and merging behaviours in parallel. Additionally, the application was
deployed publicly so attendees could also use their own laptops and devices.

\section{Participants}
Two data sources were collected:
\begin{itemize}
  \item \textbf{User Acceptance Test (UAT):} 10 in-person participants (P1--P10)
        with observer-recorded results.
  \item \textbf{Survey responses:} 13 responses collected via Microsoft Forms
        (including remote responses). The survey reported an average completion
        time of 08:46 minutes.
\end{itemize}

Survey respondents had mixed Git familiarity: 2 not familiar, 5 somewhat
familiar, 3 comfortable, and 3 very comfortable.

\section{User Acceptance Testing Tasks}
Three tasks were measured using an observer checklist to provide objective
evidence.

\subsection{Task T1: Auto-schedule Feasibility}
Scenario: at least six stops across one to two days. Evidence recorded:
completion, time-to-finish, non-overlap, and non-zero travel gaps.

\begin{table}[H]
\centering
\caption{UAT Task T1 --- Auto-schedule feasibility.}
\label{tab:uat_t1}
\begin{tabular}{@{}lcccc@{}}
\toprule
\textbf{Participant} & \textbf{Completed?} & \textbf{Time} &
\textbf{Non-overlapping?} & \textbf{Non-zero travel gaps?} \\
\midrule
P1  & Y & 3:12  & Y & Y \\
P2  & Y & 4:45  & Y & Y \\
P3  & Y & 2:58  & Y & Y \\
P4  & Y & 5:30  & Y & Y \\
P5  & Y & 3:44  & Y & Y \\
P6  & Y & 6:02  & Y & Y \\
P7  & N & 8:00* & -- & -- \\
P8  & Y & 4:11  & Y & Y \\
P9  & Y & 3:55  & Y & Y \\
P10 & Y & 4:28  & Y & Y \\
\bottomrule
\end{tabular}

\vspace{0.3em}
{\footnotesize *P7 timed out due to difficulty locating the auto-scheduler
button without a hint.}
\end{table}

Across successful completions (9/10), the mean completion time was 4:18 (median
4:11), and all successful outputs were non-overlapping with non-zero travel gaps.
This supports Success Criterion SC2.

\subsection{Task T2: Merge Conflict Resolution}
Scenario: pre-seeded repository with two branches producing at least one timing
conflict. Evidence recorded: resolved or not, time-to-resolve, and number of
hints given.

\begin{table}[H]
\centering
\caption{UAT Task T2 --- Merge conflict resolution.}
\label{tab:uat_t2}
\begin{tabular}{@{}lccc@{}}
\toprule
\textbf{Participant} & \textbf{Resolved?} & \textbf{Time} &
\textbf{Hints given} \\
\midrule
P1  & Y & 2:34  & 0 \\
P2  & Y & 4:50  & 1 \\
P3  & Y & 1:58  & 0 \\
P4  & Y & 3:22  & 0 \\
P5  & Y & 4:47  & 1 \\
P6  & N & 6:10* & 2 \\
P7  & Y & 3:05  & 1 \\
P8  & Y & 2:41  & 0 \\
P9  & Y & 4:33  & 1 \\
P10 & Y & 3:18  & 0 \\
\bottomrule
\end{tabular}

\vspace{0.3em}
{\footnotesize *P6 selected the wrong branch and ran out of time.}
\end{table}

Overall, 9/10 participants resolved the conflict; 9/10 achieved a resolution
within 5 minutes. The mean time across all participants was 3:44 (mean hints
0.6). This supports Success Criterion SC1.

\subsection{Task T3: Quick Edit + Save Without Git Explanation}
Task T3 measured whether a participant could make a small timeline/map edit and
save without requiring Git vocabulary explanation.

\begin{table}[H]
\centering
\caption{UAT Task T3 --- Quick edit and save without Git explanation.}
\label{tab:uat_t3}
\begin{tabular}{@{}lcp{0.60\linewidth}@{}}
\toprule
\textbf{Participant} & \textbf{Completed?} & \textbf{Notes} \\
\midrule
P1  & Y & Edited timeline stop duration naturally \\
P2  & Y & Used ``Save'' without prompting \\
P3  & Y & -- \\
P4  & N & Asked ``does saving here create a version?'' \\
P5  & Y & -- \\
P6  & Y & -- \\
P7  & Y & -- \\
P8  & N & Hesitated; unsure if edit was auto-saved \\
P9  & Y & -- \\
P10 & Y & -- \\
\bottomrule
\end{tabular}
\end{table}

8/10 completed without needing Git explanation. This supports the value of Easy
mode and indicates that versioning concepts can still raise questions (notably
around mental models of saving).

\section{Survey Results}
The Evaluation Day survey collected 13 responses. Key quantitative findings
included:
\begin{itemize}
  \item Git familiarity was mixed: 15\% not familiar (2/13), 38\% somewhat
        familiar (5/13), 23\% comfortable (3/13), 23\% very comfortable (3/13).
  \item Participants tried a range of features: 11 created/opened a repository,
        9 ran the auto-scheduler, 5 attempted branch/compare and merge/conflict
        resolution, and 4 used quick route mode.
\end{itemize}

Likert responses across Core Functionality and Overall Impression were strongly
positive, with responses concentrated in Agree/Strongly Agree. The Git Concept
Understanding section was also positive but showed more neutral responses,
suggesting that the conceptual model is understandable but still benefits from
guidance for some users.

The survey included the System Usability Scale (SUS). Due to the small sample
size and the visual aggregation format, this report interprets SUS qualitatively
rather than claiming a statistically robust SUS score. The overall pattern
indicates that positive usability statements were generally endorsed while
negative statements about complexity and inconsistency were generally rejected.

\subsection{Qualitative Feedback}
Free-text responses provide additional evidence of perceived value and pain
points. Representative themes include:

\textbf{Most useful/impressive:}
\begin{itemize}
  \item Auto-scheduler usefulness and edge-case handling (e.g., clearly showing
        unscheduled stops with reasons).
  \item Novelty of Git-style planning and merge conflict resolution in a non-code
        context.
  \item Educational value of algorithm visualisation.
\end{itemize}

\textbf{Most frustrating:}
\begin{itemize}
  \item Some users perceived Git concepts as requiring learning overhead.
  \item Merge conflicts can be confusing without guidance (consistent with UAT
        hint patterns).
  \item A routing limitation was reported for special transport modes (e.g.,
        ferries), reflecting dependency on external routing providers.
\end{itemize}

\textbf{Would users adopt GiTrip?} Most respondents indicated ``yes'' or
``maybe'', with collaboration and ``keeping everything in one place'' cited as
major benefits. The strongest blockers were perceived complexity for casual trips
and routing-provider constraints for niche transport types.

\section{Discussion Against Success Criteria}
Overall, the evaluation provides strong evidence that core objectives were
achieved:
\begin{itemize}
  \item The auto-scheduler reliably produced feasible plans for typical scenarios
        (SC2 met).
  \item Merge conflict resolution was achievable by most participants within the
        target time (SC1 met).
  \item Usability feedback validates the Easy mode approach while highlighting
        discoverability and conceptual learning as remaining improvement areas
        (SC3 partially met, with known limitations).
\end{itemize}

\section{Threats to Validity}
The evaluation has limitations typical of student project evaluations:
\begin{itemize}
  \item \textbf{Sample bias:} many participants were likely students with
        above-average technical literacy.
  \item \textbf{Small sample size:} limits statistical confidence, especially
        for SUS scoring.
  \item \textbf{Demonstration context:} users may perform differently in real
        trip planning over longer time horizons.
\end{itemize}

Despite these limitations, the objective task results and consistent qualitative
themes together support the conclusion that the core objectives were met.
